\section{Daten}
\label{sec:daten}

\subsection*{Datenbestand und Einheiten}

Die Primärdaten werden über den PostgREST-Dienst
\url{https://dbs.informatik.uni-halle.de/sciencedata} aus dem Schema
\code{energycharts} gelesen. Die vier in \cref{tab:views} aufgeführten Views
decken die benötigten Dimensionen Zeit, Partnerland, Erzeugungsart und Preis
ab. Der Header \code{Accept-Profile: energycharts} wählt das Schema; ein zuvor
bezogenes Bearer-Token autorisiert die Anfrage.

\begin{table}[htbp]
  \centering
  \caption{Verwendete PostgreSQL-Views und semantische Einheiten.}
  \label{tab:views}
  \begin{tabularx}{\textwidth}{@{}lXl@{}}
    \toprule
    View & Inhalt & Einheit \\
    \midrule
    \code{v\_cbpf} & bilateraler physischer Grenzfluss; positiv Import nach
      Deutschland, negativ Export aus Deutschland & GW \\
    \code{v\_cbet} & bilateraler grenzüberschreitender Handel mit derselben
      Vorzeichenkonvention & GW \\
    \code{v\_price} & Day-Ahead-Preis der Gebotszone DE--LU & EUR/MWh \\
    \code{v\_totalpower} & gesamte deutsche Nettoerzeugung einschlie{\ss}lich
      industrieller Eigenerzeugung, getrennt nach Technologien & MW \\
    \bottomrule
  \end{tabularx}
\end{table}

Die Einheit von \code{v\_totalpower} ist erklärungsbedürftig: Ihre Spaltennamen
enden auf \code{\_in\_gw}, die Werte besitzen aber die Grö{\ss}enordnung der vom
offiziellen Endpunkt \code{/total\_power} dokumentierten MW-Werte. Beispielsweise
tritt Windleistung als fünfstelliger Wert auf. Der Export interpretiert
\code{v\_totalpower} deshalb entsprechend der offiziellen
Energy-Charts-Spezifikation als MW und rechnet die Werte durch Division durch
1000 in GW um \cite{energychartsOpenapi}. Die anderen drei Grö{\ss}en werden ohne
Einheitenumrechnung übernommen.

\subsection*{Auswahl des Zeitraums}

Untersucht wird der 1. bis 31. Mai 2025 in UTC. Die Auswahl eines vollständigen
Monats vermeidet den willkürlichen Zuschnitt auf zwei besonders auffällige
Tage. 744 Stunden sind gro{\ss} genug, um Tages- und Wochenmuster, Richtungswechsel
und mehrere Preisregime zu enthalten. Gleichzeitig bleiben die 744 Spalten
einer Pixelmatrix bei einer Monatsübersicht und die 168 Spalten eines
Wochenfensters technisch handhabbar.

Die Wahl ist auch inhaltlich ergiebig. Die Datenprüfung findet 129 Stunden mit
negativem Preis, einen Minimalpreis von \(-250{,}32\mwh\) am 11. Mai um
11 Uhr UTC sowie ein Maximum von \(229{,}11\mwh\) am 19. Mai um 18 Uhr UTC.
Die erneuerbare Leistung erreicht maximal \(71{,}465\gw\). Damit enthält der
Monat hohe und niedrige erneuerbare Einspeisung sowie positive und negative
Preise, ohne dass ausschlie{\ss}lich ein vorab bekanntes Ereignis betrachtet wird.

\subsection*{Abfrage und Vorverarbeitung}

Das Skript \code{scripts/build\_postgrest\_fixture.py} setzt serverseitige
Filter für Zeit, Deutschland beziehungsweise DE--LU sowie eine Sortierung nach
Unix-Zeitstempel. Es paginiert in Seiten von höchstens 1000 Zeilen und besitzt
für jede View ein Sicherheitslimit. Dieses Vorgehen verhindert unbegrenzte
Anfragen an die gro{\ss}en Tabellen und folgt dem von PostgREST vorgesehenen
Limit-/Offset-Prinzip \cite{postgrestPagination}. Für den ausgewählten Monat
wurden 35.712 CBPF-, 35.712 CBET-, 2.976 Erzeugungs- und 744 Preiszeilen
abgerufen.

Die Rohdaten enthalten je nach View viertelstündliche oder stündliche Werte.
Um alle Variablen auf derselben Zeitachse vergleichen zu können, werden Werte
innerhalb derselben UTC-Stunde arithmetisch gemittelt. Das Mittel ist hier
sinnvoll, weil die dargestellten Grö{\ss}en Leistung beziehungsweise Preis und
nicht Energiemengen sind. Eine Summe der vier Viertelstunden-Leistungswerte
würde die Einheit verfälschen. Anschlie{\ss}end werden Erzeugungstechnologien zu
vier semantisch lesbaren Gruppen aggregiert, siehe \cref{tab:generation}.

\begin{table}[htbp]
  \centering
  \caption{Gruppierung der Erzeugungstechnologien.}
  \label{tab:generation}
  \begin{tabularx}{\textwidth}{@{}lX@{}}
    \toprule
    Gruppe & enthaltene Technologien \\
    \midrule
    Erneuerbare & Wind an Land und auf See, Solar, Biomasse, Laufwasser,
      Speicherwasser und Geothermie \\
    Kohle & Braunkohle und Steinkohle \\
    Gas & fossiles Gas \\
    Sonstige & Öl, Kohlegase, Abfall, Pumpspeicher, Kernenergie und sonstige
      Erzeugung \\
    \bottomrule
  \end{tabularx}
\end{table}

Die Zusammenfassung reduziert visuelle Komplexität, kostet aber Detail: Ein
Anstieg von Wind kann einen Rückgang von Solar innerhalb ``Erneuerbare''
überdecken. Für die Projektfrage ist die Gruppe dennoch angemessen, weil zuerst
der Zusammenhang des gesamten erneuerbaren Anteils mit Flüssen und Preis
untersucht wird. Ein späterer Drill-down auf Einzeltechnologien bleibt als
Erweiterung möglich.

\subsection*{Normalisiertes Format und Bereitstellung}

Das Exportskript schreibt \code{public/data/energy.json}. Jeder der 744
Stundeneinträge enthält Unix-Zeitstempel, deutschsprachiges Kurzlabel, vier
Erzeugungsgruppen in GW, Preis in EUR/MWh und für elf Partnerländer je einen
physischen sowie einen gehandelten Fluss in GW. Das Format orientiert sich
direkt am Elm-Domänenmodell und vermeidet, dass die Browseranwendung mehrere
gro{\ss}e Tabellen clientseitig zusammenführen muss.

Elm lädt diese Datei mit \code{Http.get}. Das ist weiterhin ein echter
HTTP-Datenpfad, zugleich aber für eine statische GitLab-Pages-Anwendung sicher:
Weder Datenbankpasswort noch Bearer-Token erscheinen im ausgelieferten
JavaScript. Der Rohdatenexport ist durch das Skript reproduzierbar.
Zugangsdaten werden ausschlie{\ss}lich beim lokalen Export verwendet und nicht
versioniert.

\subsection*{Datenqualität und Eignung}

\code{scripts/validate\_dataset.py} prüft 744 lückenlose Stundenschritte,
identische Partnerabdeckung, endliche Werte, nichtnegative Erzeugungsgruppen
und grobe GW-Plausibilitätsgrenzen. Positive und negative physische Flüsse
kommen beide vor. Der aus elf bilateralen Werten je Stunde gebildete Summenfluss
ist eine eigene explorative Kenngrö{\ss}e und wird nicht als separates
Energy-Charts-Messfeld ausgegeben.

Die Daten sind für die Forschungsfrage geeignet, weil alle benötigten
Variablen zeitlich ausgerichtet vorliegen. Grenzen bleiben: Ein Monat trägt
keine saisonale Aussage, \code{v\_totalpower} enthält industrielle
Eigenerzeugung, und aus gleichzeitiger Erzeugung und Grenzfluss lässt sich keine
physische Herkunft importierter Energie ableiten. Diese Einschränkungen werden
in Anwendung und Interpretation sichtbar gemacht.
